%!TEX root = ../template.tex
%%%%%%%%%%%%%%%%%%%%%%%%%%%%%%%%%%%%%%%%%%%%%%%%%%%%%%%%%%%%%%%%%%%%
%% abstract-pt.tex
%% NOVA thesis document file
%%
%% Abstract in Portuguese
%%%%%%%%%%%%%%%%%%%%%%%%%%%%%%%%%%%%%%%%%%%%%%%%%%%%%%%%%%%%%%%%%%%%

\ExplSyntaxOn
\cs_if_exist:NT \TitleOnHeader
  { \TitleOnHeader }
\ExplSyntaxOff

\typeout{NT FILE abstract-pt.tex}%

O paradigma da Indústria 4.0 depende intrinsecamente da virtualização dos sistemas de produção através de Digital Twins. No entanto, a criação destes modelos virtuais permanece um processo predominantemente manual, fragmentado e dispendioso, criando uma barreira significativa à sua adoção generalizada, especialmente em cenários de reconfiguração frequente. Existe, atualmente, uma descontinuidade crítica entre os dados ricos gerados durante a fase de engenharia de sistemas e os modelos semânticos necessários para a fase de operação.

Esta dissertação propõe uma framework baseada em Engenharia Conduzida por Modelos para a geração automatizada de Digital Twins dinâmicos em sistemas de manufatura. A metodologia adota AutomationML como formato de integração de dados de engenharia e Asset Administration Shell como standard alvo para a interoperabilidade e representação digital na fase operacional. Através de regras de transformação definidas, a framework processa ficheiros AutomationML para instanciar automaticamente a estrutura de Asset Administration Shell e gerar, simultaneamente, o ambiente de simulação comportamental num motor de física.

Esta abordagem valida a utilização de AutomationML como precursor essencial para a criação expedita de Asset Administration Shells, garantindo a consistência semântica entre o projeto físico e a sua contraparte digital.

\keywords{
  Digital Twin \and
  Indústria 4.0 \and
  AutomationML \and
  Asset Administration Shell \and
  Geração Automática
}